\begin{abstract}
Recovering sparse differential equations from noisy trajectories requires both estimating the underlying states and selecting terms among nearly collinear candidates. We introduce \emph{Profiled Joint Subset Search} (PJSS), which alternates joint state--coefficient estimation with subset proposals from a local quadratic model. Profiling out the states accounts for trajectory adjustment when ranking supports; exhaustive per-equation subset search explores combinations that greedy selection can miss. Promising supports are refitted with a banded solver whose cost is linear in record length, and selected by a BIC-inspired score with data-calibrated noise estimates and fixed internal defaults. On 77 in-library chaotic systems from \texttt{dysts}, PJSS recovers exact support in 89.0\% of runs at noise levels of 1--20\%, including 85.3\% on 49 systems outside our development benchmarks. In paired comparisons with full-grid DEDi, PJSS recovers 66/80 runs against 55/80 with validation selection and 61/80 with BIC selection; on 64 runs of 35 new systems, it recovers 53 against 36 and 41. Timing on three cases shows lower aggregate CPU task time and single-process GPU elapsed time. Ablations separate the larger benefit of subset proposals from the smaller additional benefit of profiling. These results concern simulated, in-library systems: measured pendulums fail the selected-support or forecast criteria, and library misspecification can produce inaccurate models without triggering the ambiguity flag.
\end{abstract}
